\documentclass[12pt,a4paper]{article}
\usepackage[utf8]{inputenc}
\usepackage[T1]{fontenc}
\usepackage[spanish]{babel}
\usepackage{geometry}
\usepackage{graphicx}
\usepackage{xcolor}
\usepackage{listings}
\usepackage{tikz}
\usepackage{booktabs}
\usepackage{array}
\usepackage{longtable}
\usepackage{hyperref}
\usepackage{fancyhdr}
\usepackage{titlesec}
\usepackage{amsmath}
\usepackage{amssymb}
\usepackage{enumitem}
\usepackage{float}

\geometry{margin=2.5cm}

\definecolor{primary}{RGB}{31,73,125}
\definecolor{secondary}{RGB}{68,114,196}
\definecolor{codebg}{RGB}{245,245,245}
\definecolor{codeframe}{RGB}{200,200,200}

\lstset{
    backgroundcolor=\color{codebg},
    frame=single,
    rulecolor=\color{codeframe},
    basicstyle=\ttfamily\small,
    keywordstyle=\color{primary}\bfseries,
    commentstyle=\color{gray}\itshape,
    stringstyle=\color{secondary},
    breaklines=true,
    showstringspaces=false,
    numbers=left,
    numberstyle=\tiny\color{gray},
}

\pagestyle{fancy}
\fancyhf{}
\fancyhead[L]{\textcolor{primary}{\textbf{Documento de Arquitectura}}}
\fancyhead[R]{\textcolor{primary}{GeoMIP — k-Particiones}}
\fancyfoot[C]{\thepage}
\fancyfoot[L]{\small Universidad de Caldas}
\fancyfoot[R]{\small 2026}

\titleformat{\section}{\Large\bfseries\color{primary}}{}{0em}{\thesection.\quad}
\titleformat{\subsection}{\large\bfseries\color{secondary}}{}{0em}{\thesubsection.\quad}
\titleformat{\subsubsection}{\normalsize\bfseries\color{black}}{}{0em}{\thesubsubsection.\quad}

\hypersetup{
    colorlinks=true,
    linkcolor=primary,
    urlcolor=secondary,
    pdftitle={Documento de Arquitectura - GeoMIP},
    pdfauthor={Jorge Ivan Garcia Torres, Camilo Botero Merchor}
}

\begin{document}

% ─────────────────────────── PORTADA ───────────────────────────
\begin{titlepage}
    \centering
    \vspace*{1cm}
    {\Huge\bfseries\color{primary} Documento de Arquitectura\\[0.5em]
    \Large Sistema GeoMIP\\[0.3em]
    \large Algoritmo k-MIP Geométrico para k-Particiones}
    \vspace{1.5cm}

    \rule{\linewidth}{2pt}
    \vspace{0.5cm}

    {\large\textbf{Proyecto de Análisis y Diseño de Algoritmos}}\\[0.3em]
    {\large Ingeniería en Sistemas y Computación}\\[0.3em]
    {\large Facultad de Ingeniería — Universidad de Caldas}
    \vspace{1cm}

    \rule{\linewidth}{0.5pt}
    \vspace{0.8cm}

    {\large\textbf{Autores:}}\\[0.5em]
    {\large Jorge Ivan Garcia Torres}\\[0.3em]
    {\large Camilo Botero Merchor}
    \vspace{1cm}

    \rule{\linewidth}{0.5pt}
    \vspace{0.5cm}
    {\large Semestre 2026-1}\\[0.2em]
    {\large\textit{feature/k-partition-benchmark}}

    \vfill
    {\small\color{gray} proyecto-analisis-kpartitions / GeoMIP}
\end{titlepage}

\tableofcontents
\newpage

% ─────────────────────────── 1. INTRODUCCIÓN ───────────────────────────
\section{Introducción}

\subsection{Propósito}
Este documento describe la arquitectura de software del sistema \textbf{GeoMIP}, una implementación del algoritmo geométrico para encontrar la partición de mínima pérdida de información integrada (k-MIP) en sistemas de N nodos. El sistema extiende la bipartición clásica de GeometricSIA hacia k-particiones arbitrarias, permitiendo evaluar sistemas de 10 hasta 25 nodos con diferentes valores de k (2, 3, 4, 5).

\subsection{Alcance}
El sistema cubre:
\begin{itemize}
    \item Generación de matrices de transición de probabilidad (TPM) para sistemas de N nodos.
    \item Ejecución del algoritmo \texttt{KPartitionGeometricSIA} sobre casos de prueba definidos en un Excel institucional.
    \item Escritura de resultados (partición, pérdida, tiempo) en el Excel del profesor.
    \item Benchmark comparativo entre Geometric y QNodes (referencia).
\end{itemize}

\subsection{Definiciones clave}
\begin{longtable}{>{\bfseries}p{3.5cm} p{10cm}}
    \toprule
    Término & Descripción \\
    \midrule
    \endhead
    TPM & Transition Probability Matrix. Matriz $2^N \times N$ donde cada fila representa un estado del sistema y cada columna la probabilidad de activación de un nodo. \\
    k-MIP & k-Partition of Minimum Information loss. Partición de k subconjuntos que minimiza la pérdida de información integrada $\Phi$. \\
    EMD & Earth Mover's Distance. Métrica de distancia entre distribuciones de probabilidad. \\
    Tabla T & Tabla de costos de transición construida mediante BFS modificado. Codifica la ``inercia causal'' entre pares de estados. \\
    $\Phi$ (Phi) & Medida de información integrada. $\Phi = 0$ indica independencia causal completa. \\
    Estado inicial & Cadena binaria de longitud N representando el estado de activación inicial del sistema. \\
    \bottomrule
\end{longtable}

% ─────────────────────────── 2. VISTA GENERAL ───────────────────────────
\section{Vista General del Sistema}

\subsection{Contexto}
El sistema GeoMIP forma parte de un proyecto de investigación académica que evalúa la eficiencia del algoritmo geométrico frente a métodos exactos (PyPhi/QNodes) para el cálculo de $\Phi$ en redes causales binarias. El contexto de uso es la plataforma de pruebas del curso de Análisis y Diseño de Algoritmos de la Universidad de Caldas.

\subsection{Diagrama de contexto}

\begin{figure}[H]
\centering
\begin{tikzpicture}[
    box/.style={rectangle, draw=primary, fill=primary!10, rounded corners, minimum width=3cm, minimum height=1cm, align=center, font=\small},
    arrow/.style={->, thick, color=primary},
    label/.style={font=\footnotesize, color=gray}
]
    \node[box] (excel_in)  at (0,4)    {Excel del Profesor\\(249 casos)};
    \node[box] (tpms)      at (5,4)    {Archivos TPM\\NxxA.csv};
    \node[box] (orquestador) at (2.5,2) {src/main.py\\(Orquestador)};
    \node[box] (algo)      at (2.5,0)  {KPartitionGeometricSIA\\(Algoritmo Core)};
    \node[box] (excel_out) at (-1,0)   {Excel Llenado\\(Resultados)};
    \node[box] (bench)     at (6,0)    {benchmark.py\\o benchmark_mejorado.py};

    \draw[arrow] (excel_in) -- (orquestador) node[midway, left, label] {casos};
    \draw[arrow] (tpms) -- (orquestador) node[midway, right, label] {TPM};
    \draw[arrow] (orquestador) -- (algo) node[midway, right, label] {condición, alcance, mecanismo, k};
    \draw[arrow] (algo) -- (orquestador) node[midway, left, label] {partición, pérdida, tiempo};
    \draw[arrow] (orquestador) -- (excel_out) node[midway, left, label] {escritura};
    \draw[arrow] (excel_out) -- (bench) node[midway, below, label] {lectura};
\end{tikzpicture}
\caption{Diagrama de contexto del sistema GeoMIP}
\end{figure}

% ─────────────────────────── 3. COMPONENTES ───────────────────────────
\section{Arquitectura de Componentes}

\subsection{Capa de Datos}

\subsubsection{Archivos TPM (data/samples/) }
Los archivos CSV representan la TPM del sistema. Su estructura es:
\begin{itemize}
    \item \textbf{Filas:} $2^N$ estados posibles del sistema.
    \item \textbf{Columnas:} $N$ nodos, valor 0 o 1.
    \item \textbf{Indexación:} little-endian. El estado \texttt{[a,b,c,...]} corresponde al índice $a + 2b + 4c + \ldots$.
\end{itemize}

\begin{longtable}{lrrr}
    \toprule
    Archivo & Nodos & Filas & Tamaño aprox. \\
    \midrule
    \endhead
    N10A.csv & 10 & 1,024     & $<1$ MB \\
    N15B.csv & 15 & 32,768    & 5 MB \\
    N20A.csv & 20 & 1,048,576 & 20 MB \\
    N22A.csv & 22 & 4,194,304 & 88 MB \\
    N25A.csv & 25 & 33,554,432 & 800 MB \\
    \bottomrule
\end{longtable}

\subsubsection{Excel del profesor (DatosPruebas2026\_1.xlsx)}
Archivo de entrada/salida principal con 5 hojas de datos y 249 casos totales.

\begin{longtable}{lrr}
    \toprule
    Hoja & Casos & Estado inicial \\
    \midrule
    \endhead
    10A-Elementos & 49 & 1000000000 (10 bits) \\
    15B-Elementos & 50 & 100000000000000 (15 bits) \\
    20A-Elementos & 50 & 10000000000000000000 (20 bits) \\
    22A-Elementos & 50 & 1000000000000000000000 (22 bits) \\
    25A-Elementos & 50 & 1000000000000000000000000 (25 bits) \\
    \bottomrule
\end{longtable}

\textbf{Mapeo de columnas de resultados:}
\begin{longtable}{ccccc}
    \toprule
    k & \multicolumn{2}{c}{QNodes (referencia)} & \multicolumn{2}{c}{Geometric (nuestro)} \\
    \cmidrule(lr){2-3}\cmidrule(lr){4-5}
      & Columnas & Contenido & Columnas & Contenido \\
    \midrule
    \endhead
    2 & D, E, F & Partición, Pérdida, Tiempo & G, H, I & Partición, Pérdida, Tiempo \\
    3 & J, K, L & Partición, Pérdida, Tiempo & M, N, O & Partición, Pérdida, Tiempo \\
    4 & P, Q, R & Partición, Pérdida, Tiempo & S, T, U & Partición, Pérdida, Tiempo \\
    5 & V, W, X & Partición, Pérdida, Tiempo & Y, Z, AA & Partición, Pérdida, Tiempo \\
    \bottomrule
\end{longtable}

\subsection{Capa de Utilidades}

\subsubsection{Manager (manager.py)}
Gestiona el sistema de N nodos. Recibe el \texttt{estado\_inicial} en binario y resuelve rutas a los archivos de muestra en \texttt{GeoMIP/data/samples/}.

\subsection{Capa de Algoritmo}

\subsubsection{KPartitionGeometricSIA (geometric\_k.py)}
Componente central del sistema. Implementa el algoritmo k-MIP geométrico.

\textbf{Flujo interno (según código):}
\begin{enumerate}
    \item \texttt{aplicar\_estrategia(condición, alcance, mecanismo, tpm)}\; prepara el subsistema vía \texttt{sia\_preparar\_subsistema}.
    \item \texttt{_construir\_tabla()}\; construye la tabla de costos \texttt{tabla\_transiciones} mediante recorrido BFS por niveles.
    \item \texttt{_candidatos\_geometricos(k)}\; genera candidatos por heurísticas basadas en los costos \texttt{T}:
    \begin{itemize}
        \item caso exacto si \texttt{n\_vars * k <= 20}
        \item caso greedy (cuantiles/rotaciones) si no.
    \end{itemize}
    \item \texttt{_evaluar\_candidato(grupos)}\; calcula la EMD real usando:
    \begin{itemize}
        \item partición por separación presente/futuro usando \texttt{n\_pres}
        \item producto tensorial con \texttt{np.kron} cuando hay múltiples distribuciones.
    \end{itemize}
    \item Selecciona el candidato con menor pérdida.
\end{enumerate}

\subsubsection{GeometricSIA (geometric.py)}
Clase base para bipartición.

\subsection{Capa de Orquestación}

\subsubsection{Orquestación Excel (src/main.py)}
Orquestador principal. Coordina la lectura del Excel, preparación de parámetros (alcance/mecanismo), ejecución del algoritmo con aislamiento por proceso y escritura de resultados.

\textbf{Características clave (implementación Method2):}
\begin{itemize}
    \item Itera los casos del Excel.
    \item Para cada caso, ejecuta la estrategia dentro de un \texttt{multiprocessing.Process}.
    \item Convierte cadenas de letras de \texttt{alcance} y \texttt{mecanismo} a binario mediante \texttt{convertir\_a\_binario}.
    \item Aplica \texttt{join(timeout=...)} y, si excede el límite, termina el proceso con \texttt{terminate()}.
    \item Carga la TPM desde \texttt{GeoMIP/data/samples/} usando las rutas resueltas por \texttt{Manager}.
\end{itemize}

\subsubsection{exec\_geometric\_k.py (ejecución por sistemas)}
Alternativa de ejecución que corre el algoritmo sobre los sistemas muestreados (N=10,15,20,22) con timeout por sistema usando multiprocessing.

\subsubsection{generar\_tpms.py}
Script auxiliar para generar los archivos TPM faltantes en \texttt{GeoMIP/data/samples/}.

\subsubsection{benchmark\_mejorado.py / benchmark.py}
Calcula métricas comparando el algoritmo geométrico contra PyPhi/QNodes (referencia), incluyendo acierto exacto, error relativo en $\Phi$, distancia de Jaccard y speedup.

% ─────────────────────────── 4. FLUJO DE DATOS ───────────────────────────
\section{Flujo de Datos}

\subsection{Conversión letras → binario}
Las cadenas del Excel contienen nodos como letras (ej. \texttt{ACEGI}). Antes de invocar el algoritmo se convierten a binario de longitud N.

\subsection{Construcción de la Tabla T}
La tabla de costos \texttt{tabla\_transiciones} se construye por BFS nivel a nivel a partir del estado inicial hasta el estado final. El factor de decrecimiento se implementa como $2^{-d_H}$ y se acumulan costos sobre caminos vecinos.

\subsection{Evaluación de candidatos (EMD)}
Para cada candidato de k-partición:
\begin{enumerate}
    \item Se separan índices de variables en presentes y futuros (según \texttt{n\_pres}).
    \item Se construye la distribución marginal por grupo usando \texttt{bipartir(...)}.
    \item Si hay más de una distribución parcial, se combina con producto tensorial (\texttt{np.kron}).
    \item Se calcula EMD (\texttt{emd\_efecto}) contra la distribución marginal del subsistema completo.
\end{enumerate}

% ─────────────────────────── 5. DECISIONES DE DISEÑO ───────────────────────────
\section{Decisiones de Diseño}

\subsection{Reducción de candidatos para k=2}
La generación ingenua de candidatos crece rápidamente con el tamaño del sistema. El enfoque documentado usa heurísticas basadas en la tabla de costos \texttt{T} del método geométrico para limitar la cantidad de candidatos evaluados.

\subsection{Timeout por proceso}
Casos pesados se ejecutan en procesos hijos con \texttt{join(timeout=...)} y terminación forzada si exceden el límite.

\subsection{Separación presente/futuro por \texttt{n\_pres}}
La evaluación de candidatos depende de distinguir índices de presentes vs futuros usando el umbral \texttt{n\_pres}; sin esta separación se producirían errores de forma al construir distribuciones.

\subsection{Estado inicial completo}
El \texttt{Manager} recibe el estado inicial del sistema completo; la proyección al subsistema activo se realiza internamente al preparar el modelo.

% ─────────────────────────── 6. DEPENDENCIAS ───────────────────────────
\section{Dependencias y Entorno}

\subsection{Dependencias Python}
\begin{longtable}{lll}
    \toprule
    Paquete & Versión mínima & Uso \\
    \midrule
    \endhead
    numpy   & 1.24 & Operaciones matriciales, TPM, EMD \\
    pandas   & 2.x & Lectura de Excel y benchmark \\
    openpyxl & 3.1 & Lectura/escritura de Excel \\
    pyinstrument & 4.x & Profiling (solo en desarrollo) \\
    \bottomrule
\end{longtable}

\subsection{Estructura de directorios (relevante)}
\begin{lstlisting}[language=bash, caption={Estructura del repositorio}]
proyecto-analisis-kpartitions/
├── GeoMIP/
│   ├── data/samples/                     # Archivos TPM (NxxA.csv)
│   └── src/Method2_Dynamic_Programming_Reformulation/
│       ├── src/main.py                   # Orquestación Excel (timeout)
│       ├── exec_geometric_k.py          # Ejecución por sistemas con multiprocessing
│       ├── benchmark.py / benchmark_mejorado.py
│       ├── controllers/strategies/
│       │   ├── manager.py
│       │   ├── geometric_k.py          # CORE
│       │   ├── geometric.py
│       │   └── q_nodes.py
│       └── generar_tpms.py
└── DocsNuevos/
    └── DatosPruebas2026_1.xlsx
\end{lstlisting}

\end{document}

